% Source: https://artofproblemsolving.com/community/c3h1901956_confused_with_domain_and_range_of_inverse_function
% Stable identity: aops-topic-1901956-q01
(a) We need $x-2\ge0$, so the domain is $[2,\infty)$. As $x$ runs over $[2,\infty)$, $\sqrt{x-2}$ takes every value in $[0,\infty)$, so the range is $[0,\infty)$.

(b) Solve $y=\sqrt{x-2}$ for $x$: $y\ge0$ and $x=y^2+2$. So
\[
\boxed{f^{-1}(x)=x^2+2\ \text{ for } x\ge0},
\]
with domain $[0,\infty)$ (the range of $f$) and range $[2,\infty)$ (the domain of $f$).

(c) An inverse function is determined by its rule and its domain. The domain of $f^{-1}$ must be the range of $f$, which is $[0,\infty)$; it does not come from the formula. On all of $\mathbb R$ the formula fails to be an inverse: $g(-1)=3$ but $f(3)=1\ne-1$, and $g$ is not even one-to-one ($g(1)=g(-1)$), so it has no inverse at all. Only after restricting $g$ to $x\ge0$ do we get $f^{-1}$.
